\chapter{Relative Value: (Alpha Book)}

\section{Regression (\texttt{pairs/stats.py})}

A pair is modeled by OLS (via \texttt{statsmodels}) of one leg on another, encapsulated in
\texttt{RegressionResults}:

\begin{itemize}
\item Stores \texttt{intercept}, \texttt{slope\_per\_step} (the \textbf{hedge ratio}), \texttt{r2}, \texttt{n\_obs}.
\item \textbf{Residual dispersion} uses $\operatorname{std}(\text{residuals}, \mathrm{ddof}=2)$ (correct degrees of freedom for a two-parameter regression) to build \textbf{confidence bands}.
\item Signals are driven by the \textbf{residual z-score} vs. configurable entry/exit bands; mean-reversion of the residual is the exit.
\end{itemize}

\section{Spread Families (Alpha Book $\rightarrow$ Spread)}

Sector PCA spreads, spread regression, treasury/policy/local/corporate spreads, swap
spreads, bond-swap spreads, and futures term/net basis. The statistical layer
(\texttt{curves/calibration/stat.py}) supplies stationarity flags and z-scores so candidates
are only flagged when the relationship is statistically stable.

\paragraph{Review focus.}
Stationarity/cointegration testing rigor, lookback choice for the regression window,
and stability of the hedge ratio through regime shifts.

\section{Alpha Portfolio $\rightarrow$ Backtest: Style-routed Spread Strategies}

The Alpha Book backtest in \texttt{web/tabs/alpha/} evaluates individual relative-value
spreads and combines the current Portfolio allocation into an indicative book
equity curve. Each spread is assigned exactly one style, mean reversion (MR) or
trend, on the first available trading day of each calendar month. The assignment
is fixed through month-end: the model does not adaptively switch engines within
the month. Candidate scores rank and allocate capital; they are not intramonth
entry confirmations.

{\small
\setlength{\LTleft}{0pt}
\setlength{\LTright}{0pt}
\setlength{\tabcolsep}{5pt}
\renewcommand{\arraystretch}{1.15}
\begin{longtable}{p{0.16\textwidth} p{0.24\textwidth} p{0.29\textwidth} p{0.25\textwidth}}
\caption{Style-routed spread strategy types.}\label{tab:pairs-strategy-types-overleaf} \\
\toprule
\textbf{Strategy type} & \textbf{Engine} & \textbf{Economic hypothesis} & \textbf{Core entry confirmation} \\
\midrule
\endfirsthead
\toprule
\textbf{Strategy type} & \textbf{Engine} & \textbf{Economic hypothesis} & \textbf{Core entry confirmation} \\
\midrule
\endhead
\midrule
\multicolumn{4}{r}{\textit{Continued on next page}} \\
\midrule
\endfoot
\bottomrule
\endlastfoot
Mean-reverting (MR) & \texttt{engine\_mr.py} & A statistically rich/cheap spread will revert toward its local equilibrium. & A 120-day rolling z-score crosses the entry band. \\
Trend & \texttt{engine\_trend.py} & A confirmed directional reversal from a local extremum will persist sufficiently to justify a directional position. & A confirmed directional-change (DC) event using the monthly frozen, volatility-calibrated absolute threshold. \\
\end{longtable}
}

\subsection{Regime Selection and Scope}
\label{sec:pairs-regime-selection}

\texttt{curves/calibration/regime.py::compute\_regime\_features} computes four rolling,
60-day indicators on first differences: Kaufman efficiency ratio, a single-scale
R/S Hurst estimate, a variance-ratio proxy, and lag-1 autocorrelation. Each casts a
trend / mean-reversion / neutral vote. A net vote of at least $+2$ selects
\texttt{trending}; at most $-2$ selects \texttt{mean\_reverting}; otherwise the result is
\texttt{uncertain}.

The classifier is invoked only at the monthly review date, using data available on
or before that date. Its result is mapped to canonical \texttt{mr} or \texttt{trend}
and held through month-end. If the result is uncertain or history is insufficient,
the prior monthly style is retained; before the first valid classification, the
spread family's static default applies. If the new style differs from an open
position's style, that position is closed once at the review observation with
\texttt{exit\_reason=monthly\_style\_change}. Regime details are retained in the review
audit trail, not used for daily routing or as a hidden entry filter.

Each monthly record contains the review date, assigned style, classifier label,
vote score/confidence, fallback reason, and, for a trend assignment, the frozen DC
threshold and its calibration inputs. This makes both the style and the timing
threshold reproducible without look-ahead bias.

\paragraph{Indicator definitions.} Let $s_t$ be the spread level, $\Delta s_t = s_t -
s_{t-1}$ its first difference, and $w=60$ observations (\texttt{DEFAULT\_REGIME\_WINDOW})
the trailing window. All four indicators are computed on the most recent $w$ values of
$\Delta s$ (or, for the efficiency ratio, on the corresponding $w$-step level change):

\begin{equation}
\begin{aligned}
\text{Efficiency ratio: } && ER_t &= \frac{\left|s_t - s_{t-w}\right|}{\displaystyle\sum_{i=t-w+1}^{t} \left|\Delta s_i\right|}, \\[4pt]
\text{Hurst exponent (single-scale R/S): } && Y_i &= \sum_{k=1}^{i}\bigl(\Delta s_k - \overline{\Delta s}\bigr),\quad
R = \max_i Y_i - \min_i Y_i,\quad S = \operatorname{sd}(\Delta s), \\
&& H_t &= \frac{\ln(R/S)}{\ln w}, \\[4pt]
\text{Variance ratio: } && w_{\mathrm{short}} &= \max\!\left(\left\lfloor w/4\right\rfloor,\,2\right),\quad
VR_t = \frac{\operatorname{Var}_{w_{\mathrm{short}}}(\Delta s)\,\bigl(w / w_{\mathrm{short}}\bigr)}{\operatorname{Var}_{w}(\Delta s)}, \\[4pt]
\text{Lag-1 autocorrelation: } && AC_t &= \operatorname{corr}\!\left(\Delta s_i,\ \Delta s_{i-1}\right) \text{ over the trailing } w \text{ observations}.
\end{aligned}
\end{equation}

Plain-language reading of each indicator:
\begin{itemize}
\item \textbf{Efficiency ratio $ER_t \in [0,1]$}: ratio of the net $w$-step displacement to
the total path length traveled to get there. Near $1$: the spread moved directly in one
direction with little back-and-forth (efficient/trending). Near $0$: large round-trip
churn relative to the net move (choppy/range-bound, favors mean reversion).
\item \textbf{Hurst exponent $H_t$}: $H_t>0.5$ indicates persistence (an up-move tends to
be followed by another up-move --- trending); $H_t=0.5$ is consistent with a random
walk (no memory); $H_t<0.5$ indicates anti-persistence (moves tend to reverse ---
mean-reverting).
\item \textbf{Variance ratio $VR_t$}: compares variance realized over the short window,
rescaled to the long window's horizon, against the actual long-window variance. $VR_t>1$
means variance is growing faster than a random walk would imply (trending/momentum);
$VR_t<1$ means it grows slower (mean-reverting, i.e.\ moves get walked back).
\item \textbf{Lag-1 autocorrelation $AC_t$}: positive autocorrelation of daily changes
indicates momentum/continuation; negative autocorrelation indicates a reversal
tendency (mean reversion).
\end{itemize}

\paragraph{Voting thresholds.} Each indicator casts one vote in $\{-1, 0, +1\}$
(mean-reverting, neutral, trending):

{\small
\setlength{\LTleft}{0pt}
\setlength{\LTright}{0pt}
\setlength{\tabcolsep}{5pt}
\renewcommand{\arraystretch}{1.15}
\begin{longtable}{p{0.24\textwidth} p{0.22\textwidth} p{0.22\textwidth} p{0.24\textwidth}}
\caption{Regime-vote thresholds per indicator.}\label{tab:pairs-regime-votes-overleaf} \\
\toprule
\textbf{Indicator} & \textbf{Trending vote ($+1$)} & \textbf{Mean-reverting vote ($-1$)} & \textbf{Otherwise} \\
\midrule
\endfirsthead
\toprule
\textbf{Indicator} & \textbf{Trending vote ($+1$)} & \textbf{Mean-reverting vote ($-1$)} & \textbf{Otherwise} \\
\midrule
\endhead
\midrule
\multicolumn{4}{r}{\textit{Continued on next page}} \\
\midrule
\endfoot
\bottomrule
\endlastfoot
Efficiency ratio $ER_t$ & $> 0.35$ & $< 0.20$ & neutral ($0$) \\
Hurst exponent $H_t$ & $> 0.55$ & $< 0.45$ & neutral ($0$) \\
Variance ratio $VR_t$ & $> 1.05$ & $< 0.95$ & neutral ($0$) \\
Lag-1 autocorrelation $AC_t$ & $> 0.05$ & $< -0.05$ & neutral ($0$) \\
\end{longtable}
}

The four votes are summed and normalized to a \texttt{regime\_score} $\in [-1, +1]$ (sum
divided by $4$); \texttt{regime\_confidence} is $\left|\texttt{regime\_score}\right|$. The
net (unnormalized) vote sum determines the label: $\ge +2$ votes selects
\texttt{trending}; $\le -2$ votes selects \texttt{mean\_reverting}; any net vote in
$\{-1, 0, +1\}$ is reported as \texttt{uncertain}. This ensemble-voting design means at
least two of the four indicators must agree in the same direction before a definite
regime is assigned; a single strong indicator outvoted by the others yields
\texttt{uncertain} rather than a confident call.

\paragraph{Not a Markov/HMM model.} Both \texttt{compute\_regime\_features} and the
directional-change trend state of \S\ref{sec:pairs-trend-model} are deterministic,
rule-based calculations on a rolling window --- there is no hidden-state inference, no
fitted transition matrix, and no probabilistic smoothing. This is a deliberate
``lightweight real-time path'' design choice (see the module docstring). Note that
\texttt{futures/backtest/regime.py} additionally defines a separate
\texttt{RegimeDetector} class built on \texttt{hmmlearn.GaussianHMM} for the futures
strategy stack; the Alpha/pairs regime classifier described here imports only the two
stateless helper functions from that module (\texttt{\_estimate\_hurst} and
\texttt{\_calculate\_variance\_ratio}) and does \textbf{not} use, fit, or depend on that
HMM class.

\subsection{Mean-reversion Model and Defaults}

For a spread level $s_t$, the MR engine uses a fixed internal lookback of 120
observations:

\begin{equation}
z_t = \frac{s_t - \operatorname{mean}_{120}(s)}{\operatorname{sd}_{120}(s)}.
\end{equation}

The engine enters long when $z_t \le -z_{entry}$ and short when
$z_t \ge z_{entry}$. A signal exit is allowed only after \texttt{min\_hold}; it occurs
when the z-score reverts inside the exit band. A $z$-score stop is always active,
including during the minimum holding period. Carry/roll and direction-aware borrow
adjustments accrue in realized PnL and can inform ranking and allocation, but do not
change an MR entry or exit decision.

{\small
\setlength{\LTleft}{0pt}
\setlength{\LTright}{0pt}
\setlength{\tabcolsep}{5pt}
\renewcommand{\arraystretch}{1.15}
\begin{longtable}{p{0.30\textwidth} p{0.16\textwidth} p{0.18\textwidth} p{0.28\textwidth}}
\caption{Mean-reversion strategy defaults.}\label{tab:pairs-mr-defaults-overleaf} \\
\toprule
\textbf{Parameter} & \textbf{Standard default} & \textbf{Tenor-spread preset} & \textbf{Role} \\
\midrule
\endfirsthead
\toprule
\textbf{Parameter} & \textbf{Standard default} & \textbf{Tenor-spread preset} & \textbf{Role} \\
\midrule
\endhead
\midrule
\multicolumn{4}{r}{\textit{Continued on next page}} \\
\midrule
\endfoot
\bottomrule
\endlastfoot
Internal z-score lookback & 120 observations & 120 observations & Local equilibrium and dispersion; not UI-configurable. \\
\texttt{entry\_z} & 2.0 & 2.5 & Absolute z-score entry threshold. \\
\texttt{exit\_z} & 0.5 & 0.25 & Reversion band used after the minimum hold. \\
\texttt{stop\_z} & 4.0 & 5.0 & Adverse z-score stop. \\
\texttt{min\_hold} & 7 calendar days & 10 calendar days & Minimum holding period for signal exits; stop remains active. \\
\end{longtable}
}

\subsection{Trend Model and Defaults}
\label{sec:pairs-trend-model}

The trend engine is event-driven. It enters long only after an \texttt{Upward Trend
Confirmed} DC event and enters short only after a \texttt{Downward Trend Confirmed}
event, when \texttt{allow\_short} is enabled. It does not use normalized momentum,
carry level, candidate score, or the current regime label as a second intramonth
entry gate.

The DC generator uses an absolute threshold in spread units:

\begin{equation}
\left|s_t-e_t\right| \ge \theta_{i,m},
\end{equation}

where $e_t$ is the active running extremum, $i$ is the spread, and $m$ is its
review month. To make the absolute event scale suitable for heterogeneous spreads,
the threshold is calibrated at the monthly review and then held fixed through
month-end:

\begin{equation}
\theta_{i,m}=\operatorname{clip}\!\left(k_{c(i)}\hat{\sigma}^{\mathrm{MAD}}_{i,m},
\theta^{\min}_{c(i)},\theta^{\max}_{c(i)}\right).
\end{equation}

Here $c(i)$ denotes the spread family and $\hat{\sigma}^{\mathrm{MAD}}$ is the
60-observation robust daily-change scale,
\begin{equation}
\hat{\sigma}^{\mathrm{MAD}}=1.4826\times
\operatorname{median}\!\left(\left|\Delta s_t-
\operatorname{median}(\Delta s)\right|\right).
\end{equation}
The family multiplier $k_{c(i)}$, noise floor $\theta^{\min}_{c(i)}$, and ceiling
$\theta^{\max}_{c(i)}$ are selected by family-level walk-forward calibration, rather
than fitted separately to every instrument. A fallback static family threshold is
used when the 60-observation history is unavailable. Optional month-to-month
smoothing limits a threshold change to $[2/3,3/2]$ of the preceding month's value.
This is still an absolute DC event rule at trade time, and is stable for inverted,
near-zero, and zero-crossing spreads.

The generator (\texttt{curves/calibration/trend.py::generate\_absolute}) is a
deterministic rule-based state machine, \textbf{not} a Markov chain or hidden Markov
model: it tracks a running local extremum and flips state only when the spread moves
away from that extremum by at least $\theta_{i,m}$. It emits an \texttt{Upward Trend
Confirmed} or \texttt{Downward Trend Confirmed} event; its state is forward-filled as
$D_t=+1$ or $D_t=-1$ until the next confirmed event.

An open position exits on an opposite confirmed DC event after \texttt{min\_hold}, or
immediately on a \textbf{volatility-scaled trailing stop}. The engine tracks the best
(most favorable) spread level reached since entry and closes the position once the
pullback reaches \texttt{trailing\_mult} $\times\hat{\sigma}^{\mathrm{MAD}}_{i,m}$.
The trailing scale is also frozen for the review month. As in MR, closed PnL combines
duration-scaled spread movement, carry accrual, and applicable borrow costs.

{\small
\setlength{\LTleft}{0pt}
\setlength{\LTright}{0pt}
\setlength{\tabcolsep}{5pt}
\renewcommand{\arraystretch}{1.15}
\begin{longtable}{p{0.24\textwidth} p{0.17\textwidth} p{0.19\textwidth} p{0.32\textwidth}}
\caption{Trend strategy parameters.}\label{tab:pairs-trend-defaults-overleaf} \\
\toprule
\textbf{Parameter} & \textbf{Standard default} & \textbf{Tenor-spread preset} & \textbf{Role} \\
\midrule
\endfirsthead
\toprule
\textbf{Parameter} & \textbf{Standard default} & \textbf{Tenor-spread preset} & \textbf{Role} \\
\midrule
\endhead
\midrule
\multicolumn{4}{r}{\textit{Continued on next page}} \\
\midrule
\endfoot
\bottomrule
\endlastfoot
DC calibration lookback & 60 observations & 60 observations & Point-in-time robust daily-change scale at the monthly review. \\
DC threshold $\theta_{i,m}$ & $\operatorname{clip}(k_c\hat{\sigma}^{\mathrm{MAD}},\theta_c^{\min},\theta_c^{\max})$ & Same family rule & Absolute threshold, fixed through the review month. \\
Family multiplier $k_c$ & Walk-forward calibrated & Walk-forward calibrated & Controls DC sensitivity; estimated by family, never per-instrument in sample. \\
Threshold floor / cap & Family configured & Family configured & Prevents quote-noise churn and post-shock signal starvation. \\
Threshold smoothing & $[2/3,3/2]$ prior month & $[2/3,3/2]$ prior month & Optional bound on monthly threshold changes. \\
\texttt{trailing\_mult} & 1.5 & 2.0 & Multiple of the frozen monthly robust scale. \\
\texttt{min\_hold} & 7 calendar days & 10 calendar days & Minimum hold before an opposite-DC signal exit. \\
\texttt{allow\_short} & enabled & enabled & Allows short-spread entries. \\
\end{longtable}
}

\subsection{Model Comparison and Deliberate Exclusions}

The selected design separates three decisions: monthly style selection, monthly
threshold calibration, and event-level trade timing. Table
\ref{tab:pairs-model-comparison-overleaf} records the alternatives considered and
why they are not used as the production entry model.

{
\small
\setlength{\LTleft}{0pt}
\setlength{\LTright}{0pt}
\setlength{\tabcolsep}{5pt}
\renewcommand{\arraystretch}{1.15}
\begin{longtable}{p{0.24\textwidth} p{0.27\textwidth} p{0.38\textwidth}}
\caption{Trend-model comparison and production decision.}\label{tab:pairs-model-comparison-overleaf} \\
\toprule
\textbf{Alternative} & \textbf{Potential benefit} & \textbf{Why it is not the production entry rule} \\
\midrule
\endfirsthead
\toprule
\textbf{Alternative} & \textbf{Potential benefit} & \textbf{Why it is not the production entry rule} \\
\midrule
\endhead
\midrule
\multicolumn{3}{r}{\textit{Continued on next page}} \\
\midrule
\endfoot
\bottomrule
\endlastfoot
Single global absolute DC threshold & Simple and transparent. & Spread levels, volatility, liquidity, and quote noise differ materially by instrument. One global threshold creates excessive flips in quiet series and too few events in volatile series. The selected model instead uses a monthly frozen, family-calibrated absolute threshold. \\
Relative DC threshold & Automatically changes with the level of a positive series. & It is unstable or undefined near zero and changes directional interpretation for negative/inverted yield spreads. The production rule compares an absolute movement in spread units. \\
Normalized momentum entry & Measures fixed-window move strength relative to volatility. & It is lagging and path-insensitive: a new DC reversal can be vetoed because its trailing window still contains the preceding move. It also needs an arbitrary lookback and threshold. It remains a diagnostic and may inform threshold research, but does not confirm, veto, reverse, or size a trade. \\
DC plus normalized-momentum confirmation & May reduce some weak DC events. & It introduces a second correlated gate, delays entry, and obscures the causal reason for a trade. The monthly style assignment is the broad filter; the confirmed DC event is the single timing rule. \\
Carry-level or candidate-score gate & May favour positive expected carry and highly ranked candidates. & Carry is realized in PnL and candidate scores determine ranking/allocation, but neither establishes a trend reversal. Using either as an entry veto mixes portfolio selection with trade timing and creates a hidden signal filter. \\
Daily adaptive MR/trend routing & Responds quickly to changing market conditions. & Daily switching increases turnover and makes historical attribution and auditability weak. It can switch an open position without a separate trade event. The selected model permits a style change only at the next monthly review. \\
HMM / Markov regime model & Can estimate latent states and transition probabilities. & It adds fitted-state, transition, and retraining risk for limited per-spread histories, while reducing transparency. It is not required for the baseline. HMM research remains separate from this deterministic Alpha model. \\
\end{longtable}
}

The chosen model is therefore not a claim that momentum or HMM methods cannot add
value. It is a parsimonious baseline with a single observable entry event and a
complete point-in-time audit trail. Any alternative must demonstrate incremental
out-of-sample net performance, stable turnover, and comparable interpretability
before it replaces this baseline.

For yield-based spread families (\texttt{TenorSpread}, \texttt{SwapSpread},
\texttt{TBondCurve}/\texttt{CBondCurve}, \texttt{TBondSwap}/\texttt{CBondSwap},
\texttt{NetBasis}, \texttt{FuturesSwap}, \texttt{PCASpread}, \texttt{BinarySpread}),
the callback negates the raw spread ($-s_t$) before passing it to either engine, so
that the engine's native $\mathrm{LONG}$ (position $=+1$) already means ``expect the
raw yield-based spread to fall/narrow,'' and native $\mathrm{SHORT}$ already means
``expect the raw spread to rise/widen.'' Prices, $z$-scores, and score series are
sign-restored for display afterwards. This convention must be retained when
reviewing direction, carry, and threshold settings.

\paragraph{Direction-label bug (identified and fixed 2026-07-16).} The individual
backtest callback previously applied an \emph{additional} $\mathrm{LONG}$/$\mathrm{SHORT}$
swap to the trend engine's output only (not to the MR engine's output) after the
sign-restoration step above. Because the engine's direction was already correct
under the $-s_t$ convention, this extra swap double-flipped every trend trade's
displayed direction for yield-based spreads: genuine $\mathrm{SHORT}$ signals (e.g.
during a sustained \texttt{CGB-10s30s} widening move) were shown as $\mathrm{LONG}$
and vice versa, which could be misread as an absence of short signals. PnL figures
were unaffected (they are computed from the underlying position sign, not the
display label). The redundant swap has been removed from
\texttt{web/tabs/alpha/callbacks/backtest\_tab.py}, so trend trades now use the same
(correct) direction convention as MR trades. The trend score-history chart was also
updated to display the same sign convention explicitly.

\paragraph{Candidate \texttt{trend\_state} key bug (identified and fixed 2026-07-16).}
\texttt{curves/refreshers/alpha\_scoring.py} populated the candidate table's
\texttt{trend\_state} column by reading the key \texttt{"state"} from
\texttt{compute\_trend\_signal}'s return dictionary, but that function returns the
value under the key \texttt{"trend\_state"}. The lookup therefore fell back to its
default of \texttt{0.0} for effectively every candidate, silently suppressing the
directional-change state shown in the candidate scan even when a trend was
confirmed. The key has been corrected.

\subsection{Portfolio Aggregation and Measurement Limitations}

Portfolio mode loads the persisted current Alpha snapshot, normalizes the positive
candidate weights across instruments with available overlapping history, and sums
their weighted daily mark-to-market PnL in basis points. It forwards the MR
entry/exit/stop/minimum-hold controls, but trend assets use the trend engine's
function defaults rather than the individual-screen trend controls. The resulting
book is a useful \textbf{current-book sanity view}, not a fully parameterized historical
simulation.

The portfolio screen accepts initial capital and a transaction-cost input (default
100 MM and 0.5 bp), but these values are parsed and are not applied to PnL in the
current callback. Likewise, the trend path in portfolio mode does not receive its
spread-type/borrow-cost arguments. Reported portfolio return is consequently
weighted cumulative PnL in bp, rather than capital-normalized net return. The
individual engines report a trade-PnL Sharpe,
$\operatorname{mean}(\mathrm{pnl})/\operatorname{sd}(\mathrm{pnl})\times\sqrt{\min(N_{\mathrm{trades}},20)}$,
whereas the portfolio screen reports a daily-PnL Sharpe annualized by
$\sqrt{252}$. These are not comparable with each other or directly with the
platform-wide return-based convention in \S1.

\subsection{Validation and Implementation Requirements}

Before increasing risk or selecting family parameters, the implementation must meet
the following validation requirements:

\begin{enumerate}
\item \textbf{Make performance measurement investable first.} Apply instrument- and
direction-specific bid/ask, fees, financing, borrow, and conservative next-bar
execution assumptions on every entry, exit, reversal, and portfolio rebalance;
then make capital and the configured transaction cost operational. Report gross
and net results, turnover, capacity/DV01, and return-based daily Sharpe using one
shared convention.

\item \textbf{Validate with a genuinely walk-forward design.} Freeze the monthly style,
robust scale, threshold multiplier, floor/cap, and trailing-stop setting using only
information available at each review date; reserve a final untouched period. Use
purged, embargoed, or anchored walk-forward folds where overlapping holding periods
make ordinary splits optimistic. Compare against always-MR, always-trend, no-trade,
and the rejected alternatives in Table \ref{tab:pairs-model-comparison-overleaf}.

\item \textbf{Calibrate the monthly classifier and threshold by family.} The
variance-ratio proxy and single-scale 60-day R/S Hurst estimate have finite-sample
limitations. Calibrate classifier votes and the DC multiplier/floor/cap to
instrument-family null distributions and walk-forward results. The review output
must always report confidence, fallback reason, robust scale, and final absolute
threshold.

\item \textbf{Enforce one causal DC implementation.} All trend paths must call the
absolute DC generator with the review month's frozen threshold. Test positive,
inverted, near-zero, and zero-crossing spreads; verify no entry occurs before the
first confirmed event; and verify that normalized momentum and carry cannot gate an
entry or an opposite-DC exit.

\item \textbf{Estimate parameters by spread family, not individual in-sample fit.}
Estimate an out-of-sample OU half-life/stationarity diagnostic per instrument; trade
MR only when its half-life fits the intended 1-week to 1-month holding horizon. For
trend, select the DC multiplier, floor/cap, robust-scale lookback, and trailing-stop
multiple by conservative family-level grids, scoring stability across folds rather
than the single best historical Sharpe.

\item \textbf{Improve portfolio construction after signals are credible.} Preserve the
historically available allocation at each rebalance; volatility-target gross DV01;
cap key-rate/issuer/leg and correlated-family exposure; net shared legs; and apply
drawdown, liquidity, and event-risk overlays. Add cross-sectional rich/cheap
signals within a spread family to reduce common level-factor risk that a collection
of independent time-series signals cannot diversify away.
\end{enumerate}
